\documentclass[11pt]{article}
\usepackage[a4paper,margin=18mm]{geometry}
\usepackage[no-math]{fontspec}
\setmainfont{Latin Modern Roman}
\usepackage{amsmath,amssymb,amsthm,xfrac,xspace,hyperref,graphicx,comment}
\excludecomment{DefTralics}
\providecommand{\C}{}
\providecommand{\bibliofont}{\small}
\newcommand{\ie}{i.e.,\xspace}
\newcommand{\eg}{e.g.,\xspace}
\newcommand{\etc}{etc.\xspace}
\newcommand{\cf}{cf.\xspace}
\newcommand{\resp}{resp.\xspace}
\newcommand{\smaller}{\small}
\newcommand{\Subsection}{\subsection}
\newcommand{\Subsubsection}{\subsubsection}
\newcommand{\skpt}{\smallskip}
\emergencystretch=3em
\input{author-preamble.tex}
\begin{document}
\begin{center}{\Large Stable item 1935}\end{center}
\paragraph{Source and attribution.}
Jean-René Chazottes, Pierre Collet, Servet Martínez, Sylvie Méléard, \emph{Quasi-stationary distributions and resilience: what to get from a sample?}, Journal de l’École polytechnique — Mathématiques 7 (2020), publisher version, DOI 10.5802/jep.132. \href{https://jep.centre-mersenne.org/articles/10.5802/jep.132/}{Publisher article}. Authors' text: \href{https://creativecommons.org/licenses/by/4.0/}{CC BY 4.0}.
\paragraph{Editorial problem boundary.}
Prove only the one-species (\(d=1\)) instance of Theorem 1.1 below, with all printed hypotheses (H.1)--(H.8) and the stated large-radius condition on the moment estimates. For a fixed lag \(\tau\ge0\), the process is started in its quasi-stationary law and the killed-process lag covariance satisfies the printed \(\Sigma^K(\tau)=e^{\tau M^*}\Sigma^K(0)+O(\sqrt K)\) asymptotic. The full moment and centring core and generator/Taylor/variation-of-constants proof remain editable. No multitype, statistical-estimator or growing-lag claim is selected.
\paragraph{Difficulty (editorial): H2.}
Exponential quasi-stationary concentration and a Lyapunov recursion yield centered moments at Gaussian scale. Generator differentiation and Taylor remainders then identify the population correlation semigroup with the linearized deterministic return rate to order square-root population size, despite eventual absorption.
\paragraph{Editorial transcription note.}
Original author language, mathematics and ordering are preserved. Publisher front matter is replaced by this independent article wrapper. Bibliography is recovered from matching cached publisher HTML, including its TeX formulas. No new proof or mathematical repair is supplied. Surrounding original results are retained for conservative context and are not extra selected corpus claims.
\clearpage
\input{author-body.tex}
\begin{thebibliography}{99}
\bibitem{patrick} P. Billingsley - Convergence of probability measures , Wiley Series in Probability and Statistics , John Wiley \& Sons, Inc., New York, 1999
\bibitem{ccm1} J.-R. Chazottes, P. Collet \& S. Méléard - “Sharp asymptotics for the quasi-stationary distribution of birth-and-death processes” , Probab. Theory Relat. Fields 164 (2016) no. 1-2, p. 285-332
\bibitem{ccm2} J.-R. Chazottes, P. Collet \& S. Méléard - “On time scales and quasi-stationary distributions for multitype birth-and-death processes” , Ann. Inst. H. Poincaré Probab. Statist. 55 (2019) no. 4, p. 2249-2294
\bibitem{EK} S. N. Ethier \& T. G. Kurtz - Markov processes. Characterization and convergence , Wiley Series in Probability and Statistics , John Wiley \& Sons, Inc., New York, 1986
\bibitem{resilience-book} - Foundations of ecological resilience (L. H. Gunderson, C. R. Allen \& C. S. Holling, eds.), Island Press, 2012
\bibitem{Ikeda} N. Ikeda \& S. Watanabe - Stochastic differential equations and diffusion processes , vol. 24 , North-Holland Publishing Co., Amsterdam, Tokyo, 1989
\bibitem{kubo} R. Kubo - “The fluctuation-dissipation theorem” , Rep. Progr. Phys. 29 (1966), p. 255-284
\bibitem{sim} V. Simoncini - “Computational methods for linear matrix equations” , SIAM Rev. 58 (2016) no. 3, p. 377-441
\end{thebibliography}
\end{document}
